% Corte mecánico íntegro: parte_iv.tex:220--336.
% Residencia material del ensamblaje sucesor de 102 capítulos.
% Residencia de realización: capítulo 58.
\chapter{Electrón central y lector helicoidal de la materia}

\section{Unicidad del centro material bajo inversión celular}

El atlas material tiene dominio \(\mathcal C_{81}=\{1,\ldots,9\}^2\) e
involución
\[
 \rho_c(r,c)=(10-r,10-c).
\]
Su único punto fijo es
\[
 (r,c)=(5,5).
\]
Las otras \(80\) celdas forman \(40\) pares. El electrón central es la
única multisección puntual; las demás especies internas son fibras, órbitas
o multisecciones. La teoría espiral no cuestiona esta conclusión.

\section{Modo trítico común}

En la carta canónica APP \(1,\ldots,9\), el vector de valores del selector
trítico de fase operativa es
\[
 m_{\ph}=(1,-1,0)^3.
\]
Para la tabla conjunta tridimensional,
\[
 C^{(3)}m_{\ph}=(C^{(3)})^{\!*}m_{\ph}
 =-\frac12m_{\ph}.
\]
Por tanto
\[
 s=\frac12,\qquad
 1-s^2=\frac34=\frac{90}{120},
\]
\begin{equation}
 \|[P_{\Sigma,3},P_{\Pi,3}]\|
 =\frac{\sqrt3}{4}.
 \label{espiral:eq:modo-electron-app}
\end{equation}
El modo, el defecto orientado \(90\to120\) y el prefactor electrónico son
publicaciones coordinadas de una misma cadena APP--TRIT--TPK, cuya composición
efectiva proporciona su entrelazamiento interno.

\section{Ledger material}

Cada celda del atlas porta la firma
\[
 \mathcal R(r,c)=
 (n_A,s_C,k_{\Delta_4},\nu_{120},\nu_{270},\tau_{12})_{r,c}.
\]
La imagen contiene \(56\) familias de ruta; el refinamiento por familia y
nivel contiene \(13\) multisecciones. Las cifras rectoras son
\[
 81/56/13/324,
\]
con el sector nulo y las torres globales. El lector radial sólo puede añadirse
mediante una firma compatible, por ejemplo
\[
 \Sigma_{\mathrm{esp}}(X)=
 (p_X,a_X,[\Lambda_X,C_X],
 [\kappa_{9,X}],\mu_X,\varepsilon_{{\mathrm{tw}},X}),
\]
que sea constante donde corresponda y natural bajo refinamiento.

\section{Cuadrado de compatibilidad material}

Sea \(\mathcal E_{\mathrm{atlas}}\) el espacio de secciones persistentes y
\(\mathfrak S_{\mathrm{esp}}\) el lector espiral. La compatibilidad material
se expresa mediante el cuadrado
\[
\begin{array}{ccc}
\mathcal E_{\mathrm{atlas}}
 & \xrightarrow{\ \mathfrak S_{\mathrm{esp}}\ }
 & \mathcal D_{\rad}^{\enr}\\[1.0em]
{\scriptstyle (K_D,K_\Omega)}\big\downarrow
 && \big\downarrow{\scriptstyle \ev_{\mathrm{dim}}^{\mathrm{esp}}}\\[1.0em]
\operatorname{Spec}(K_D,K_\Omega)
 & \xrightarrow{\ \mathfrak R_{\mathrm{mat}}\ }
 & \mathcal M_{\mathrm{hel}}
\end{array}
\]
El electrón \((5,5)\) proporciona el anclaje central, pero no selecciona
retrospectivamente \(p,\Lambda,C\). Para las demás especies se requiere
constancia sobre multisecciones, compatibilidad con inversión celular,
\(K_D/K_\Omega\), CKM/PMNS, color y representación de espín.

\section{Espín y helicidad}

No se identifica el espín \(1/2\) con una vuelta de una hélice. La
representación fermiónica satisface
\[
 U(2\pi)\psi_e=-\psi_e,
 \qquad
 U(4\pi)\psi_e=\psi_e.
\]
La hélice nonádica registra fase y memoria; la representación de espín vive
en otro codominio. Un entrelazador material válido debe conservar ambas
leyes, no reemplazar una por la otra.

\begin{figure}[htbp]
\centering
\begin{tikzpicture}[scale=.63]
\foreach \i in {1,...,9}{
 \foreach \j in {1,...,9}{
  \draw[HMTLine,fill=HMTLight] (\i,\j) rectangle +(1,1);
 }
}
\fill[HMTNavy] (5.15,5.15) rectangle (5.85,5.85);
\draw[HMTRed,very thick] (5.5,5.5) circle (1.05);
\draw[->,HMTRed,thick] (5.5,5.5)--(8.9,8.9);
\node[right,HMTNavy,font=\sffamily\bfseries] at (8.9,8.9) {único punto fijo \((5,5)\)};
\node[below,HMTGray,font=\small\sffamily] at (5.5,.45) {atlas material \(9\times9\)};
\end{tikzpicture}
\caption{El electrón central es un resultado del atlas; la correspondencia espiral es un lector posterior.}
\label{espiral:fig:atlas-electron}
\end{figure}
